\subsection{Two infinite endpoints and one exact exponent}
\label{rw:infinity:two-infinite}

When both ray endpoints tend to infinity, the lower cones and upper
barrier of Section~\ref{rw:finite:geometry} give the limiting conditions
\[
 \bar B\ge0,\qquad |\bar a|\le2\sqrt{\bar B},\qquad
 \bar c^2+4\bar m\le0.
\]
The case $\bar c=\bar m=0$ is reserved for the second scale in
Section~10. In this subsection $\bar m<0$. These inequalities describe
the limiting point. The neighborhoods below are open in the coefficients and retain perturbations on both
sides of the limiting boundaries.

For $\bar B>B_*>0$, the two lower passages have exactly the same
power exponent at every positive scale. By
Proposition~\ref{rw:lower-preparation} and
\eqref{rw:transitions:equal-exponents},
\begin{equation}\label{rw:infinity:horizontal-powers}
 D_\pm(t)=K_\pm t^{h/\epsilon}e^{P_\pm(t)/\epsilon},\qquad
 h=\frac{1-\epsilon\bar B}{\bar B},\quad P_\pm(0)=0.
\end{equation}
The selected jets of the scaled amplitudes $P_\pm$ are bounded;
a regular radial coordinate change need not bound $P_\pm/\epsilon$.
All input-independent constants, including slow actions, remain in
$K_\pm>0$. With the two weak references fixed, the complete central
equation is
\begin{equation}\label{rw:infinity:closing}
 F=D_+(t)-J(t)-\rho\nu H_c,\qquad J=D_-\circ f,
 \quad H_c=H_++H_->h_*>0,
\end{equation}
where $s_+=-\rho\nu H_+$, $s_-=\rho\nu H_-$ and $H_\pm>0$.
The upper connector has $f'>0$. Both deviation derivatives are positive,
so $K=\epsilon\log(J'/D_+')$ satisfies
$F'=D_+'(1-e^{K/\epsilon})$. Impose every slow-path and
regular-connection condition on the comparison domain, and central
positivity on actual returns. A zero of \eqref{rw:infinity:closing} has
positive common height and positive core arcs by the argument of
Section~\ref{rw:finite:three-saddle}. With the collar and stated
passage conditions it gives an actual return.

On a strict focus neighborhood the upper angular velocity is bounded
away from zero between fixed horizontal sections. Its bounded-time
radial equation gives $f=tk$, with positive bounded $k,k^{-1},f'$
and inverse derivative, in all selected jets. Apply
Lemma~\ref{rw:joint-jet-budget}, with the common exponent
$h$, to obtain \eqref{rw:transitions:common-power}. Its coefficient-only
constant is $C=\epsilon\log(K_-/K_+)$.
All logarithm arguments have positive margins on a fixed small
section. The cancellation uses equality of the exponents for the same
field, before differentiation. On each endpoint block,
Lemma~\ref{rw:joint-jet-budget} bounds the input-dependent terms
through relative connector units. Input-independent slow actions stay
in $C$; their possibly divergent coefficient derivatives are outside
these finite-jet estimates.

\subsection{Strict focus: values and slopes after center division}

\begin{proposition}\label{rw:infinity:strict-focus}
At a limiting point with
$\bar B>0$, $\bar a^2<4\bar B$, $\bar m<0$ and
$\bar c^2+4\bar m<0$, one common neighborhood of the two infinite
endpoints contains at most two isolated counted cycles, with
multiplicity. The statement includes both center families and their
intersection.
\end{proposition}
\begin{proof}
Use the actual common input from
Lemma~\ref{rw:transitions:common-input}, so that its limit is $t=1/p$.
On compact positive input intervals, the upper regular convergence
and lower rate integrals give, on $\nu=0$,
\[
 K(t)\longrightarrow-\mathcal A(1/t),\qquad
 \mathcal A(p)=A_-(\mathcal U_0(p))-A_+(p),\qquad
 A_\pm(P)=\int_0^P\frac{x\,\dd x}{1\pm\bar ax+\bar Bx^2}.
\]
This convergence holds with the selected coefficient derivatives.
It extends through $t=0$: reserve one further input derivative in
\eqref{rw:transitions:common-power}, compare on a short interval to a fixed
positive input, and use the common Lipschitz bound before sending that
input to zero. On this strict domain $C$ and its selected coefficient
jets are bounded; no slow gap is tending to zero in this assertion.

To count near the two center families, we need the value and slope
of the limiting action after division by their generators.
Put $M=-\bar m>0$, $\Omega=\sqrt{4M-\bar c^2}$ and
$\Upsilon=e^{\pi\bar c/\Omega}$. The upper linear solution in
Section~\ref{rw:finite:geometry} has
\[
 S(T)=e^{\bar cT/2}\frac{\sin(\Omega T/2)}{\Omega/2},\qquad
 I(T)=\int_0^TS(u)\,\dd u,\qquad pS(T)=I(T).
\]
At $T_0=2\pi/\Omega$, $S=0$, $S'=-\Upsilon$,
$S''=-\bar c\Upsilon$, and $I=(1+\Upsilon)/M$.
Expanding $S(T)=I(T)/p$ by the implicit function theorem yields
\begin{equation}\label{rw:infinity:upper-series}
 \mathcal U_0(p)=\Upsilon p-\frac{\bar c(1+\Upsilon)}M
          +\frac{\Upsilon-\Upsilon^{-1}}{2Mp}+O(p^{-2}).
\end{equation}
Substitution in $q=(S^2-IS')/S$ gives this expansion.
The lower integrands give
\[
 A_\pm(p)=\frac{\log p}{\bar B}+C_\pm
 \mathbin{\pm}\frac{\bar a}{\bar B^2p}
 +\frac{\bar B-\bar a^2}{2\bar B^3p^2}+O(p^{-3}),
 \quad C_--C_+=\frac{\pi\bar a}{\bar B\sqrt{4\bar B-\bar a^2}}.
\]
The constant follows by separating the logarithmic derivative of
$1\pm\bar ax+\bar Bx^2$ from its reciprocal. Consequently
$\mathcal A(1/t)=J_0+J_1t+O(t^2)$, where
\begin{equation}\label{rw:infinity:focus-jets}
\begin{aligned}
 J_0&=\frac\pi{\bar B}
 \left(\frac{\bar c}{\sqrt{4M-\bar c^2}}
       +\frac{\bar a}{\sqrt{4\bar B-\bar a^2}}\right),\\
 J_1&=\frac{1+\Upsilon^{-1}}{\bar B^2}
               \left(-\bar a-\frac{\bar B\bar c}M\right).
\end{aligned}
\end{equation}
These vanish simultaneously only on $\bar a=\bar c=0$ or
$M=\bar B,\bar a=-\bar c$. Indeed $J_0=0$ fixes the opposite
signs and gives $M\bar a^2=\bar B\bar c^2$; inserting the resulting
$\bar a=-\bar c\sqrt{\bar B/M}$ in $J_1$ proves the assertion.

The real period domains in Section~\ref{rw:centers} contain every
represented input at either center. On a smaller product neighborhood,
Proposition~\ref{rw:generator-division} therefore gives
\[
 K=\tau K_\tau+\sigma_1K_\sigma+\nu K_\nu,
 \qquad \tau=\bar a+\bar c,
 \quad \sigma_1=\bar c(\bar B-M).
\]
The independent dividing coordinates are
$(\tau,\bar c,\bar B-M,\nu)$; no inverse for $\sigma_1$ is used.
The value and slope distinguish the two transverse action columns.
Differentiating \eqref{rw:infinity:focus-jets} by the successive
differences of Proposition~\ref{rw:generator-division} gives the
value--slope determinants
\[
 -\frac{2\pi(1+\Upsilon^{-1})}
 {\bar B^3(4\bar B-\bar c^2)^{3/2}},\qquad
 -\frac{\pi}{\bar B^3M(\sqrt{\bar B}+\sqrt M)}
\]
on the second and reversible center, respectively. They agree at
the intersection and have uniform nonzero margins on compact blocks.
Reserve the two coefficient
orders for this division, in addition to the three input orders and
the one derivative spent in the cancellation formula. The primitive
order can be fixed before the common sections by Section~4.

Here $\epsilon\log K_\pm$ are bounded. Choose $t_0$ so that half
of $t^{h/\epsilon}$ absorbs these constants; the other half absorbs
the finite differentiated powers and logarithms, exactly as in
\eqref{rw:finite:three-flat}. Center division therefore gives the
absolute weight required by Lemma~\ref{rw:centers:relative-forcing},
with $I=|\tau|+|\sigma_1|$. Every fibre with a closing zero satisfies
$|\nu|\le\eta(\epsilon)I$, $\eta\to0$.
If $I=|\tau|+|\sigma_1|>0$, the rank just computed yields
\[
 |K(0)|+|K'(0)|\ge cI,\qquad \|K\|_{C^2}\le CI.
\]
On a common short interval these estimates force either a nonzero
value throughout or a first derivative of one strict sign. Thus $K$
has at most one zero.
Off the centers the same alternative follows directly from $(J_0,J_1)$.
If $I=0$, forcing gives the center identity. Finally the displayed
formula for $F'$ gives the claimed two zeros, with multiplicity.
\end{proof}

\subsection{Fixed slow strips and their relative limits}

As a lower double root or an upper repeated root is approached, transit
times diverge and the strict-focus bounds on the action constants no
longer apply. Fixed sections separate these contractions from the
relative terms that will be differentiated.

\begin{lemma}\label{rw:infinity:fixed-strips}
Near an incoming lower double root $\bar a=2\sqrt{\bar B}$,
choose fixed lower phase sections on opposite sides of its minimum.
The passage exists precisely when its exact scalar factor $f_L$ is
positive between them. Its map, about the selected weak reference, is
$L_Lu e^{r_L(u)}$, with bounded selected relative jets and
\[
 \epsilon\log L_L=-A_L+O(\epsilon),\qquad A_L\longrightarrow+\infty.
\]
Near an upper repeated root $\bar c=2\gamma>0$,
$\bar m=-\gamma^2$, fix the upper slow sections $v=\pm v_R$.
That passage exists precisely for $\mu_U>0$ and has
\begin{equation}\label{rw:infinity:upper-slow-factor}
 G(u)=e^{-A_U}u e^{r_U(u)},\qquad
 A_U=\frac{2\Xi}{b\sqrt{\mu_U}}\arctan\frac{v_R}{\sqrt{\mu_U}},
\end{equation}
where $b=1-\epsilon\bar B$,
$\mu_U=-[(\bar c+\nu)^2+4b\bar m]/(4b^2)$ and
$\Xi=(\bar c+\nu)(1-2\epsilon\bar B)/(2b)$.
The complete upper connector is
$f(t)=\alpha t e^{\psi(t)}$, with
$\alpha=e^{-A_U}k_0$, $k_0>c$, and bounded selected jets of
$\psi$ and of the absolute factor $\alpha$.
These bounds do not bound coefficient derivatives of $A_L$ or $A_U$.
\end{lemma}
\begin{proof}
Apply the exact phases and relative integral estimates of Section~4.
On the lower strip, the normal rate has a positive margin and the
scalar velocity is $\epsilon f_L$. At its limiting minimum
$r_0$, positivity and coefficient closeness give
$f_L(w)\le\delta+C(w-r_0)^2$. For fixed sections $R<r_0<Q$,
\[
 A_L=\int_R^Q\frac{\beta_L(w,0)}{f_L(w)}\,\dd w
 \ge\frac{2\beta_*}{\sqrt{C\delta}}
       \arctan\left(c_1\sqrt{C/\delta}\right)\longrightarrow\infty,
 \quad c_1=\min(r_0-R,Q-r_0).
\]
The weak-branch reference may arrive at the outer section with
nonzero normal value. Subtract its solution before applying the
relative integral, so the difference still has the contracting normal
factor. Regular central composition puts $-A_L+O(1)$ into
$\epsilon\log K_-$ and leaves $P_-$ bounded in the selected
input-dependent jets. Reflection gives the corresponding assertion
for $K_+$.

The upper exact equations are $v'=\mu_U+v^2$ and
$R'=-R\beta_U(v,R)$, with
$\beta_U(v,0)=(\Xi-v)/b$.
Integration between $\pm v_R$ gives
\eqref{rw:infinity:upper-slow-factor}. Both endpoints are fixed and
separated, so the lost-exit extension of the relative-core proposition
applies. On adjoining sections $v_R\le|v|\le v_{\rm out}$,
$\mu_U+v^2\ge v_R^2/2$ after shrinking coefficients. The regular
maps $\Phi_{\rm in},\Phi_{\rm out}$ have bounded positive units.
The absolute factor $e^{-A_U}$ and the relative integral have the
selected extensions from Proposition~\ref{rw:relative-core}.
Lemma~\ref{rw:app:flat-composition} gives the complete connector
and its asserted jets.
\end{proof}

\begin{lemma}\label{rw:infinity:repeated-tail}
After subtraction of the input-independent contractions in the
preceding lemma, the relative input derivatives converge to their
repeated-root limits, uniformly even when the transit times diverge.
On compact intervals $p>1/\gamma$ the incoming upper logarithmic
rate has limit $-p/(p-1/\gamma)^2$.
\end{lemma}
\begin{proof}
Apply Proposition~\ref{rw:relative-core} to the fixed strips.
The upper strip has $q=1$, fixed transverse exits and quadratic time
jets; the lower strip has $q=\epsilon$ and $T\ge c/\epsilon$.
Their auxiliary flow jets converge on bounded time intervals, so
Corollary~\ref{rw:transitions:relative-continuity} gives the
relative differentiated limits even as the transit times diverge.

Its limiting rate is explicit. For the repeated affine field set
$V=Y+\gamma^{-2}$, $Z=X+\gamma Y-\gamma^{-1}$ and $r=Z/V$.
Then $V'=\gamma V+Z$, $Z'=\gamma Z$, $r'=-r^2$, and
$\log(Vr)-\gamma/r$ is constant. At $(p,0)$ it is
$\log(p-P_0)-P_0/(p-P_0)$, where $P_0=1/\gamma$.
The incoming relative logarithm has the negative of this entry value;
its derivative is $-p/(p-P_0)^2$. Fixed regular connectors preserve
this rate on the original entry coordinate. The preceding tail
argument supplies its perturbed limit. Reflection uses its own forward
contracting chart.
\end{proof}

\subsection{Domination and competing slow actions}

\begin{proposition}\label{rw:infinity:slow-domination}
At an incoming lower double root with a strict upper focus, the local
bound is one cycle, necessarily simple and attracting. Its reflected
outgoing configuration is repelling and has no closing zero for
$\ell\ge0$. At a positive repeated upper root with a regular outgoing
lower path, the bound is again one simple attracting cycle; it is absent
for $\ell\le0$. This includes an additional incoming lower double root.
The reflected upper configuration reverses both stability and the
excluded sign of $\ell$.
\end{proposition}
\begin{proof}
Every fixed slow strip must pass; outside its scalar passing condition
there is no represented orbit. In the first case
\eqref{rw:transitions:common-power} is $K=-A_L+O(1)$, uniformly in input
and transit time. In the second it is
$K=-hA_U-A_L+O(1)$, where $A_L=0$ if the incoming lower path
is regular. Shrink the neighborhoods until the divergent negative terms
absorb the bounded remainder. This uses no ratio between the actions and
gives $K<-1$ and $F'>0$ throughout the comparison interval. At a
zero the first-return multiplier in central height is
$J'/D_+'<1$, giving simplicity and attraction.
The deviations both tend to zero as $t\downarrow0$, so in the
second case $D_+>J$ excludes $\ell\le0$ in
\eqref{rw:infinity:closing}. Reflection gives $D_+<J$ in the
outgoing-only lower case and excludes $\ell\ge0$. These inequalities
are valid on the extended comparison interval; they do not assert a
positive central return at $t=0$.
\end{proof}

\begin{proposition}\label{rw:infinity:opposing}
Suppose the limiting data are
\[
 \bar B_*>0,\qquad \bar a_*=-2\sqrt{\bar B_*},\qquad
 \bar c_*=2\gamma>0,\quad\bar m_*=-\gamma^2,\qquad
 \bar B_*\ne\gamma^2.
\]
A common neighborhood of the two infinite endpoints contains at most
two isolated counted cycles, with multiplicity. The same holds after
reflection and for either sign of the line-breaking parameter.
\end{proposition}
\begin{proof}
Here the lower and upper contractions enter $K$ with opposite signs.
Both fixed slow strips must pass. The lower contraction belongs to
$K_+$, so cancellation gives
\[
 K(0)=A_L-hA_U+O(1),\qquad
 \|K-K(0)\|_{C^2([0,t_0])}\le C.
\]
The value $K(0)$ can be arbitrary because the two gaps vanish
independently. Keep it undifferentiated in coefficients. To
bound the nonconstant part, reserve at least three input derivatives
of the primitive relative maps before the derivative spent in the
critical logarithm, and use the same actual outgoing coordinate
before either slow strip.

Put $r_0=1/\sqrt{\bar B_*}$ and $P_0=1/\gamma$.
The regular outgoing entry integral has logarithmic rate
$-p/[\bar B_*(p-r_0)^2]$. Combining this with the upper rate in
Lemma~\ref{rw:infinity:repeated-tail}, and with the limiting horizontal
factor $h=1/\bar B_*$, gives on compact subsets of
$p>\max(r_0,P_0)$
\[
 K_p\longrightarrow\frac p{\bar B_*}
      [(p-r_0)^{-2}-(p-P_0)^{-2}].
\]
In the reciprocal input this is
\begin{equation}\label{rw:infinity:opposing-slope}
 K_t\longrightarrow
 \frac{(P_0-r_0)[2-(P_0+r_0)t]}
 {\bar B_*(1-r_0t)^2(1-P_0t)^2}.
\end{equation}
Compare with a fixed positive input and use the uniform second
derivative bound to extend this convergence to $t=0$. Since
$P_0\ne r_0$, the right side has one strict sign on a common short interval. Thus $K$
has at most one zero regardless of its constant value, and $F$ at
most two. The excluded equality is exactly the joint center treated
in Section~\ref{rw:joint-section}.
\end{proof}

\subsection{The cubic horizontal boundary}

\begin{proposition}\label{rw:infinity:cubic}
At $\bar B_* =\bar a_*=0$, $\bar m_*<0$,
$\bar c_*^2+4\bar m_*\le0$, the two-infinite-endpoint local bound
is three isolated cycles, with multiplicity. The open neighborhoods
include both signs of $\bar B$, both signs of $\nu$, and a repeated
upper boundary when present.
\end{proposition}
\begin{proof}
At the cubic boundary the power representation is singular.
Use instead the undivided phases of
Proposition~\ref{rw:lower-preparation}, written as
\[
 x'=\epsilon g_\pm(x),\quad g_\pm(x)=xF_\pm(x),\qquad
 Z'=-Z\beta_\pm(x,Z).
\]
The exact coefficients satisfy
\[
 F_\pm(0)=\bar B,\quad F_\pm'(0)=\pm\bar a,
 \quad \beta_\pm(0,0)=1-\epsilon\bar B,
\]
and on $\nu=0$,
$F_\pm(x)=\bar B\pm\bar ax+(1+\epsilon\bar m)x^2$.
Fix a small exit $x_R>0$ with $F_\pm(x_R)>x_R^2/2$ and a
positive fixed fast input. An entry $x<x_R/2$ passes precisely when
$F_\pm>0$ on $[x,x_R]$. This is the positive scalar component
adjoining $x_R$; it is one interval. Lost exits have infinite time,
so the relative and absolute extensions of Section~4 apply here.
No extension across a collision of entry and exit is needed.

After regular continuation to the central section set
\[
 N_\pm=\beta_\pm(x,0)+\epsilon g_\pm\widetilde j_\pm',\qquad
 A_\pm=N_\pm-\epsilon g_\pm'
                    +\epsilon g_\pm N_\pm'/N_\pm.
\]
The exact logarithmic rate is
\begin{equation}\label{rw:infinity:cubic-rate}
 D_\pm'>0,\qquad
 \frac{D_\pm''}{D_\pm'}=\frac{A_\pm(x)}{\epsilon xF_\pm(x)},
 \quad A_\pm=1+O(\epsilon),\quad A_\pm(0)=1-2\epsilon\bar B.
\end{equation}
The last two assertions hold with the reserved jets on the auxiliary
product box. Normal contraction absorbs the finite inverse scale powers
from the central connection, giving the absolute bounds required for
one coefficient division and the subsequent forcing estimate.

For a strict upper focus, $f=tk$ and its limiting map is
$f_0(t)=1/\mathcal U_0(1/t)$. At a positive upper repeated root,
Lemma~\ref{rw:infinity:fixed-strips} gives instead
$f=\alpha te^\psi$, $\alpha\to0$. Its proof only requires bounded
$\bar B$, so remains valid here. In both cases
\[
 L=tf'/f,\qquad R_f=f''/f',\qquad L(0)=1
\]
have bounded input jets. The compositions depend on $f$ and $\psi$,
not on $\alpha^{-1}$. The two scalar passing intervals at $t$ and
$f(t)$ intersect in one interval, since $f'>0$; at the repeated
root impose also $\mu_U>0$. The signed central comparison may again
be larger than the actual return domain.

The connector curvature enters when differentiating the complete
critical function. Keeping this term gives the exact identity
\begin{equation}\label{rw:infinity:cubic-numerator}
\begin{split}
 \mathcal P&=tF_+(t)F_-(f(t))K'\\
 &=F_+(t)A_-(f(t))L-F_-(f(t))A_+(t)
        +\epsilon tF_+(t)F_-(f(t))R_f.
\end{split}
\end{equation}
The exact endpoint values give $\mathcal P(0)=0$, allowing one
division by the radial input on the auxiliary box. Define
$Q=\mathcal P/t=\int_0^1\mathcal P'(ut)\,\dd u$ on the whole
auxiliary box. Reserve primitive total order eight: the rate and
curvature formulas spend at most two orders, and this division one.
The remaining five supply three input and two coefficient orders for
$Q$. On the passing interval $Q=F_+F_-K'$, with positive multiplier.

If $\bar c_*\ne0$, reflect to $\bar c_*>0$. At zero lower
coefficients and zero scale,
$Q=t[L-(f/t)^2]$ and
\[
 Q'(0)=1-\alpha_0^2>0,\qquad
 \alpha_0=e^{-\pi\bar c_*/\sqrt{-4\bar m_*-\bar c_*^2}}
\]
in the strict focus case. In the repeated case the same exact
numerator gives
\[
 |Q'(0)-(1-\alpha^2)|\le C(|\bar B|+|\bar a|+\epsilon),
 \qquad \|Q''\|\le C.
\]
Since $\alpha\to0$, a common small interval has $Q'>c>0$.
It follows that $Q$ has at most one zero, $K$ at most two, and
$F$ at most three.

It remains to include $\bar c_*=0$, in a strict focus neighborhood.
Choose the reflected phase and graph extensions symmetrically. On
$\bar a=\bar c=\nu=0$ the connector is an increasing involution
and hence the identity; the two relative maps and the extended
numerator coincide. This is an identity on the common auxiliary
product box, even if a contracted coefficient point loses a horizontal
exit. Thus Hadamard division gives
\[
 Q=\bar a Q_a+\bar cQ_c+\nu Q_\nu.
\]
The same division of the absolute deviations and
Lemma~\ref{rw:centers:relative-forcing} give, at a closing zero,
$|\nu|\le\eta(\epsilon)(|\bar a|+|\bar c|)$, $\eta\to0$.

Compute the transverse rank from the undivided limiting numerator,
then transfer it to the actual divided columns.
Put $M=-\bar m$ and expand the reciprocal of
\eqref{rw:infinity:upper-series} as
$f_0(t)=\alpha_0t+f_2t^2+f_3t^3+O(t^4)$. Then
\[
 f_2=\frac{\bar c\alpha_0(1+\alpha_0)}M,\qquad
 f_3=\frac{\bar c^2\alpha_0(1+\alpha_0)^2}{M^2}
                -\frac{\alpha_0-\alpha_0^3}{2M}.
\]
At zero scale the value and slope of $Q$ are
\[
 Q(0)=\bar B f_2/\alpha_0+\bar a(1+\alpha_0),
\]
\[
 Q'(0)=\bar B(2f_3/\alpha_0-f_2^2/\alpha_0^2)
       +\bar af_2(1+\alpha_0^{-1})+1-\alpha_0^2.
\]
Differentiating in $(\bar a,\bar c)$ on the reversible face gives
\[
 \begin{pmatrix}2&2\bar B/M\\
 0&\pi(M-\bar B)/M^{3/2}\end{pmatrix}.
\]
Its determinant has a fixed nonzero margin near $\bar B=0$.
Mixed convergence of the undivided extended numerator transfers this
rank to the actual columns. With
$I=|\bar a|+|\bar c|$, forcing and the reserved bounds give
$|Q(0)|+|Q'(0)|\ge cI$, $\|Q\|_{C^2}\le CI$.
The value-or-slope alternative on a fixed short interval gives at most
one zero of $Q$, followed by the same two Rolle steps. If $I=0$,
forcing gives the reversible identity. Apply the finite-selection
convention of Section~\ref{rw:counting}.
\end{proof}
